\section*{Resumo}

O problema de Minimizar Vértices \textit{d-Branch} em uma Árvore Geradora (d-MBV) consiste em, dado um grafo conexo, não direcionado e sem pesos, identificar uma árvore geradora que minimize o número de vértices com grau superior a $d$, para $d \ge 2$. Esse desafio tem aplicação direta na alocação otimizada de \textit{switches} WDM em redes ópticas voltadas à transmissão \textit{multicast}. O trabalho atual investiga a relação entre a densidade do grafo e a quantidade de vértices com grau maior que $d$ nas soluções, propondo uma metodologia que combina a análise da densidade com a centralidade \textit{PageRank} para guiar um novo algoritmo construtivo. Essa abordagem híbrida é então integrada a diferentes meta-heurísticas, como o \textit{Multi-Start} e o GRASP (\textit{Greedy Randomized Adaptive Search Procedure}), com o objetivo de explorar a busca de soluções de forma mais eficiente. Os experimentos computacionais, realizados sobre instâncias conhecidas do problema, evidenciam a eficácia da proposta, que alcança soluções de alta qualidade.

\textbf{Palavras-chave:} Otimização Combinatória. Meta-heurística. Multi-Start. GRASP. Problema d-MBV.